% answers_short.tex —— 2012 年·天津卷·文科 · 简短答案速查
% 由 tools/make_short_answers.py 自 answers.tex 的 \gksolhead/\gkanswer 提取，
% 并与 2012/2012_数学真题_天津文_答案要点.md 逐题比对后生成，勿手改（改源 MD 后重跑）。

\begin{gknotice}
\textbf{使用说明}：本节只给最终答案便于快速核对，完整推导见解析版。答案为编写者独立推导，\textbf{非官方标准答案}。
\end{gknotice}

\gkpart{选择题}
\begin{center}\begin{tabular}{ccccccccc}
\toprule
题号 & 1 & 2 & 3 & 4 & 5 & 6 & 7 & 8 \\
\midrule
答案 & C & B & C & A & A & B & D & B \\
\bottomrule
\end{tabular}\end{center}

\gkpart{填空题}
\begin{itemize}[leftmargin=2.6em]
  \item[\textbf{9.}] $-3$。
  \item[\textbf{10.}] $30$。
  \item[\textbf{11.}] $a=1$，$b=2$。
  \item[\textbf{12.}] $3$。
  \item[\textbf{13.}] $\dfrac43$。
  \item[\textbf{14.}] $(0,1)$。
\end{itemize}

\gkpart{解答题}
\begin{itemize}[leftmargin=2.6em]
  \item[\textbf{15.}] (1) 小学 $3$ 所，中学 $2$ 所，大学 $1$ 所；(2) ① 见解析（共 $15$ 种）；② $\dfrac15$
  \item[\textbf{16.}] (1) $\sin C=\dfrac{\sqrt7}4$，$b=1$；(2) $\dfrac{\sqrt{21}-3}8$
  \item[\textbf{17.}] (1) $2$；(2) 证明见解析；(3) $\dfrac{\sqrt{39}}{13}$
  \item[\textbf{18.}] (1) $a_n=3n-1$，$b_n=2^n$；(2) 证明见解析
  \item[\textbf{19.}] (1) $\dfrac{\sqrt6}4$；(2) $\pm\sqrt5$
  \item[\textbf{20.}] (1) 增区间 $(-\infty,-1)$，$(a,+\infty)$；减区间 $(-1,a)$；(2) $\left(0,\dfrac13\right)$；(3) $\dfrac43$
\end{itemize}
